En base a los datos obtenidos de las mediciones, se puede concluir que la complejidad asintótica predice bien la forma de las curvas, pero no nos cuenta toda la historia. En el ordenamiento los cuatro algoritmos crecieron como $n \log n$ en $D_7$, pero aun así \texttt{std::sort} fue 4 veces más rápido que Patience Sort a $n = 10^7$ (576 vs 2304 ms), por la diferencia de constante y de localidad de memoria: por más que las cotas sean similares, la implementación y el uso de recursos son distintos.
Con las matrices ocurre lo mismo, pero de forma más extrema. Strassen creció con factor $7.0$ al duplicar $n$, exactamente lo que predice $O(n^{\log_2 7})$, y sin embargo a $n=1024$ tardó $130$ s frente a $1$ s de naïve y usó $50$ MB de memoria auxiliar vs $4$ MB, porque el costo de copiar submatrices en cada llamada pesa más que las operaciones que ahorra. La diferencia se va acortando al crecer $n$ ($296\times \to 130\times$), pero el cruce queda fuera de nuestra medición: para matrices de tamaño pequeño y mediano, Strassen no es competitivo frente al algoritmo ingenuo.
